% =============================================================
% §6 Discussion
% =============================================================
\section{Discussion}
\label{sec:discuss}

\subsection{Implications for AIGC Scheduling Research}
\label{sec:implications}

Our experiments support three implications that extend beyond
the specific schedulers compared in \S\ref{sec:eval}.

\textbf{Simulator physics dominate algorithmic sophistication.}
The largest effect in our study---the
$-83\%$ \slo{} collapse from disabling continuous batching
(Table~\ref{tab:ablation})---is an order of magnitude larger
than any reward-design or state-engineering choice we
measured. Research effort is better invested in faithful
physical simulation than in network architectures or reward
formulations.

\textbf{Pretraining is the cheapest essential investment.}
Disabling pretrain cost $38\%$ \slo{} ($p<0.001$), yet only
$20$ episodes ($\sim$10--15~s wall-clock) suffice. Pretraining
matters not for sample efficiency but for \textbf{escaping
locally optimal regions}: earlier configurations collapsed
into an ``always-cloud'' attractor until extended pretrain was
paired with the cloud-overuse penalty.

\textbf{Don't over-attribute gains to reward shaping.} No
individual AIGC reward component (or even all jointly)
significantly improves
\slo{} once continuous batching and pretraining are in place,
and the weight-sensitivity study (\S\ref{sec:weight-sens})
shows the \emph{relative weighting} of all seven terms can
shift from latency-first to AIGC-first without measurable
effect. Prior works claiming gains from specific reward
components may have overlooked this
\textbf{substitutability} within an expressive PPO backbone:
a positive single-reward result is \emph{suspicious} without
the kind of joint ablation in Table~\ref{tab:ablation}.

Collectively, these implications suggest the AIGC scheduling
agenda should reorient from ``engineering better rewards'' to
\textbf{``engineering better simulators and training
procedures.''}

\subsection{Limitations and Threats to Validity}
\label{sec:limitations}

\textbf{Why simulation, and its fidelity.} We evaluate in
simulation deliberately. Our central claims are
\emph{comparative and statistical}---a Pareto-dominance
relation that must hold across 30 configurations and
$N{=}30$ seeds each (900 runs). Reproducing this sweep on
real hardware would require simultaneous access to A100-,
T4-, and Jetson-class nodes under controlled, repeatable
arrivals---precisely what a simulator isolates: identical
workloads and seeds across all eleven schedulers, no confound
from co-tenancy, thermal throttling, or driver drift. We calibrate the
simulator's constants (cold-load times, batching overheads,
the memory-bandwidth floor) to published vLLM and
TensorRT-LLM measurements so that absolute magnitudes are
realistic, not merely relative.

That said, the simulator makes two simplifications. It models
continuous batching as \emph{admission-time} batching rather
than vLLM's iteration-level dynamic re-batching (possibly
underestimating true batching benefit by 10--20\% in some
regimes), and treats the memory-bandwidth floor as constant
rather than batch-size-dependent. We therefore read absolute
ablation magnitudes (e.g., $-83\%$ from \texttt{no\_batching})
as our simulator's behavior rather than vLLM-precise
predictions; the \emph{ordering} and \emph{significance} of
effects---which carry our claims---are robust to these
constants. \textbf{Validating the Pareto relation on a
physical cloud--edge testbed with a production serving engine
is the most important item of future work}, and the open
simulator we release is designed to make that comparison
direct: its \texttt{scheduler.schedule()} interface is the
same one a real dispatcher would implement.

\textbf{Workload and hardware coverage.} Our workload uses
synthetic log-normal lengths calibrated to Azure 2023---not
replayed production traces or bursty arrivals---and covers
A100/T4/Jetson tiers but not H100 or custom accelerators. The
qualitative physics are architecture-agnostic; constants need
re-calibration for new hardware.

\textbf{Statistical power and algorithm scope.} Our $N=30$
runs detect medium-to-large effects (Cohen's $d \ge 0.5$);
the ablation nulls mean
``no detectable difference $>5\%$ in \slo,'' not ``no
difference.'' PPO, A3C-R2N2, and the GNN variant all
exhibited similar AIGC substitutability, but
offline/model-based RL or imitation learning might respond
differently.

\subsection{Practical Deployment Guidance}
\label{sec:guidance}

Our results offer four pieces of practical
guidance. First, \textbf{capacity planning matters more than
scheduler choice}: 70B requests overflowing the single
cloud server (\S\ref{sec:workload}) pin all schedulers to
${\approx}10\%$ \slo---\emph{provision
compute first; optimize scheduling second}. Second,
\textbf{simple load balancers suffice in the moderate-load
regime}: LeastLoaded
and ShortestQueue achieve $78\%$ of RL's \slo{}
with zero training cost. Third, \textbf{reach for RL when
energy efficiency, tail latency, or dispatch-time budget
matter}: RL attains the lowest
\etok{} in all 30 configurations (margins up to 9\% under
load), its tail TTFT
(P95 25.2~s vs.\ 28--41~s) suits
interactive serving, and its sub-millisecond policy
inference avoids the $\approx$70~ms-per-dispatch search cost
of NSGA-II-style schedulers in the request
path. Fourth, \textbf{be skeptical of
reward-engineering claims}: without joint ablation against a
pretrained baseline, gains may reflect general RL backbone
improvements rather than the claimed mechanism.
